% GENERATED FILE. Edit canonical lessons, then run npm run generate:books.
% canonical-source-hash: fnv1a64:66950e017965382f
% canonical-lessons: LA-C343-captivus, LA-C343-dedo, LA-C343-deditio, LA-C343-foedus, LA-C343-triumphus

\chapter{Triumphus}
\label{ch:la-triumphus}
\hlchaptermodality{343}

\begin{chapteropening}
\textbf{By the end of this chapter:} \emph{I can tell how a war ended: a prisoner, to surrender, a surrender, a treaty and the triumph.}

Complete the last lesson of chapter 343: I can tell how a war ended: a prisoner, to surrender, a surrender, a treaty and the triumph.
\end{chapteropening}

\section[captīvus, captīvī]{captīvus, captīvī — a prisoner, a captive}

\label{lesson:LA-C343-captivus}



\begin{quote}
Before the new one: say the Latin for to lay waste, then the Latin for a defeat.
\end{quote}

\subsection*{You'll want to know: captīvus, captīvī}
\textbf{captīvus, captīvī} — \textquotedblleft{}a prisoner, a captive\textquotedblright{}.

\begin{quote}
Captīvī post currum imperātōris ambulābant.
\end{quote}

\begin{quote}
\emph{The prisoners were walking behind the general's chariot.}
\end{quote}

Masculine, from \emph{capiō} (taught). Also an adjective, \emph{captīvus, -a, -um}. English \emph{captive}.

\begin{cousinweb}
\textbf{captīvus} is built on \emph{captus}, the participle of \emph{capiō}, \textquotedblleft{}to take\textquotedblright{} (Lewis \& Short). English \emph{captive} comes from it, and \emph{caitiff} through French.
\end{cousinweb}

\begin{grammarlens}[title={What you've built}]
One more word for telling, explaining and describing at B1.
\end{grammarlens}

\subsection*{Your turn}
\begin{itemize}
\raggedright
  \item \emph{Say it:} \emph{captīvus, captīvī}
  \item \emph{Say it:} \emph{captīvus, captīvī}, once more
  \item \emph{Say it:} a sentence of your own with \emph{captīvus, captīvī}
  \item \emph{Recall:} say the Latin for to lay waste, then the Latin for a defeat, then say \emph{captīvus, captīvī} again
\end{itemize}

\begin{tcolorbox}[breakable,colback=teal!4,colframe=teal!35!black,title={Before you move on}]
What does captīvus, captīvī mean? (\textquotedblleft{}A prisoner, a captive\textquotedblright{}.) Say it once more.
\end{tcolorbox}
\section[dēdō, dēdere]{dēdō, dēdere — to give up, to surrender}

\label{lesson:LA-C343-dedo}



\begin{quote}
Before the new one: say the Latin for a defeat, then the Latin for a prisoner.
\end{quote}

\subsection*{You'll want to know: dēdō, dēdere}
\textbf{dēdō, dēdere} — \textquotedblleft{}to give up, to surrender\textquotedblright{}.

\begin{quote}
Cīvēs sē et urbem imperātōrī dēdidērunt.
\end{quote}

\begin{quote}
\emph{The citizens surrendered themselves and the city to the general.}
\end{quote}

\emph{dēdō, dēdere, dēdidī, dēditum}: \emph{dē} + \emph{dō}. \emph{Sē dēdere}: to surrender.

\begin{cousinweb}
\textbf{dēdō} is built on \emph{dō}, \textquotedblleft{}to give\textquotedblright{}, with \emph{dē}, \textquotedblleft{}away\textquotedblright{} (EtymDB): to give away from oneself.
\end{cousinweb}

\begin{grammarlens}[title={What you've built}]
One more word for telling, explaining and describing at B1.
\end{grammarlens}

\subsection*{Your turn}
\begin{itemize}
\raggedright
  \item \emph{Say it:} \emph{dēdō, dēdere}
  \item \emph{Say it:} \emph{dēdō, dēdere}, once more
  \item \emph{Say it:} a sentence of your own with \emph{dēdō, dēdere}
  \item \emph{Recall:} say the Latin for a defeat, then the Latin for a prisoner, then say \emph{dēdō, dēdere} again
\end{itemize}

\begin{tcolorbox}[breakable,colback=teal!4,colframe=teal!35!black,title={Before you move on}]
What does dēdō, dēdere mean? (\textquotedblleft{}To give up, to surrender\textquotedblright{}.) Say it once more.
\end{tcolorbox}
\section[dēditiō, dēditiōnis]{dēditiō, dēditiōnis — a surrender}

\label{lesson:LA-C343-deditio}



\begin{quote}
Before the new one: say the Latin for a prisoner, then the Latin for to give up.
\end{quote}

\subsection*{You'll want to know: dēditiō, dēditiōnis}
\textbf{dēditiō, dēditiōnis} — \textquotedblleft{}a surrender\textquotedblright{}.

\begin{quote}
Post longam obsidiōnem dēditiō facta est.
\end{quote}

\begin{quote}
\emph{After a long siege a surrender took place.}
\end{quote}

Feminine, from \emph{dēdō}. In a \emph{dēditiō}, the defeated put everything into the winner's hands.

\begin{cousinweb}
\textbf{dēditiō} is built on \emph{dēdō}, \textquotedblleft{}to give up, to surrender\textquotedblright{} (EtymDB), with the ending \emph{-tiō}.
\end{cousinweb}

\begin{grammarlens}[title={What you've built}]
One more word for telling, explaining and describing at B1.
\end{grammarlens}

\subsection*{Your turn}
\begin{itemize}
\raggedright
  \item \emph{Say it:} \emph{dēditiō, dēditiōnis}
  \item \emph{Say it:} \emph{dēditiō, dēditiōnis}, once more
  \item \emph{Say it:} a sentence of your own with \emph{dēditiō, dēditiōnis}
  \item \emph{Recall:} say the Latin for a prisoner, then the Latin for to give up, then say \emph{dēditiō, dēditiōnis} again
\end{itemize}

\begin{tcolorbox}[breakable,colback=teal!4,colframe=teal!35!black,title={Before you move on}]
What does dēditiō, dēditiōnis mean? (\textquotedblleft{}A surrender\textquotedblright{}.) Say it once more.
\end{tcolorbox}
\section[foedus, foederis]{foedus, foederis — a treaty, an alliance}

\label{lesson:LA-C343-foedus}



\begin{quote}
Before the new one: say the Latin for to give up, then the Latin for a surrender.
\end{quote}

\subsection*{You'll want to know: foedus, foederis}
\textbf{foedus, foederis} — \textquotedblleft{}a treaty, an alliance\textquotedblright{}.

\begin{quote}
Duae cīvitātēs foedus fēcērunt.
\end{quote}

\begin{quote}
\emph{The two states made a treaty.}
\end{quote}

Neuter, 3rd declension. Not the adjective \emph{foedus, -a, -um}, 'foul'. English \emph{federal}.

\begin{cousinweb}
\textbf{foedus} is inherited from the Proto-Indo-European root \emph{*b\textsuperscript{h}eyd\textsuperscript{h}-} (EtymDB), which also gives \emph{fidēs}, \textquotedblleft{}trust\textquotedblright{} (taught); Lewis \& Short link it with \emph{fīdō}, \textquotedblleft{}to trust\textquotedblright{}. English \emph{federal} comes from it.
\end{cousinweb}

\begin{grammarlens}[title={What you've built}]
One more word for telling, explaining and describing at B1.
\end{grammarlens}

\subsection*{Your turn}
\begin{itemize}
\raggedright
  \item \emph{Say it:} \emph{foedus, foederis}
  \item \emph{Say it:} \emph{foedus, foederis}, once more
  \item \emph{Say it:} a sentence of your own with \emph{foedus, foederis}
  \item \emph{Recall:} say the Latin for to give up, then the Latin for a surrender, then say \emph{foedus, foederis} again
\end{itemize}

\begin{tcolorbox}[breakable,colback=teal!4,colframe=teal!35!black,title={Before you move on}]
What does foedus, foederis mean? (\textquotedblleft{}A treaty, an alliance\textquotedblright{}.) Say it once more.
\end{tcolorbox}
\section[triumphus, triumphī]{triumphus, triumphī — a triumph, a victory procession}

\label{lesson:LA-C343-triumphus}



\begin{quote}
Before the new one: say the Latin for a surrender, then the Latin for a treaty.
\end{quote}

\subsection*{You'll want to know: triumphus, triumphī}
\textbf{triumphus, triumphī} — \textquotedblleft{}a triumph, a victory procession\textquotedblright{}.

\begin{quote}
Imperātor in triumphō per urbem vectus est.
\end{quote}

\begin{quote}
\emph{The general rode through the city in a triumph.}
\end{quote}

Masculine. The winner rode in a chariot through the city, and his soldiers sang rude songs about him.

\begin{cousinweb}
\textbf{triumphus} was taken, EtymDB says, from Greek \emph{thriambos}, \textquotedblleft{}a hymn to Dionysus\textquotedblright{}, probably through Etruscan; old Latin wrote \emph{triumpus}. English \emph{triumph} comes from it.
\end{cousinweb}

\begin{grammarlens}[title={What you've built}]
One more word for telling, explaining and describing at B1.
\end{grammarlens}

\subsection*{Your turn}
\begin{itemize}
\raggedright
  \item \emph{Say it:} \emph{triumphus, triumphī}
  \item \emph{Say it:} \emph{triumphus, triumphī}, once more
  \item \emph{Say it:} a sentence of your own with \emph{triumphus, triumphī}
  \item \emph{Recall:} say the Latin for a surrender, then the Latin for a treaty, then say \emph{triumphus, triumphī} again
\end{itemize}

\begin{tcolorbox}[breakable,colback=teal!4,colframe=teal!35!black,title={Before you move on}]
What does triumphus, triumphī mean? (\textquotedblleft{}A triumph, a victory procession\textquotedblright{}.) Say it once more.
\end{tcolorbox}
